\documentclass[11pt]{article}
\usepackage[T1]{fontenc}
\usepackage[utf8]{inputenc}
\usepackage[a4paper,margin=1in]{geometry}
\usepackage{lmodern}
\usepackage{amsmath,amssymb}
\usepackage{graphicx}
\usepackage{listings}
\usepackage{booktabs}
\usepackage{caption}
\usepackage{subcaption}
\usepackage{hyperref}

\title{Quick Instructions for Beamer Slides}
\author{}
\date{}

\begin{document}
\maketitle

\section*{Create a Slide}
Basic frame:
\begin{verbatim}
\section{Section Title}
\begin{frame}{Slide Title}
  Your bullet points / content here.
\end{frame}
\end{verbatim}

Long slides (auto-break):
\begin{verbatim}
\begin{frame}[allowframebreaks]{Long Title}
  ...
\end{frame}
\end{verbatim}

Code on a slide (use \texttt{[fragile]}):
\begin{verbatim}
\begin{frame}[fragile]{Code Example}
\begin{lstlisting}[language=Python]
for i in range(5):
    print(i)
\end{lstlisting}
\end{frame}
\end{verbatim}

\section*{Lists}
\begin{verbatim}
\begin{itemize}
  \item First
  \item Second
\end{itemize}

\begin{enumerate}
  \item Step 1
  \item Step 2
\end{enumerate}
\end{verbatim}

\section*{Math}
Inline: \verb|$a^2+b^2=c^2$| \\
Display:
\begin{verbatim}
\[
  E = mc^2
\]
\end{verbatim}
Aligned:
\begin{verbatim}
\begin{align*}
  \nabla\cdot \mathbf{E} &= \frac{\rho}{\varepsilon_0},\\
  \nabla\times \mathbf{E} &= -\frac{\partial \mathbf{B}}{\partial t}.
\end{align*}
\end{verbatim}

\section*{Images}
Single image (centered):
\begin{verbatim}
\begin{frame}{Figure}
  \centering
  \includegraphics[width=0.7\linewidth]{figures/my_plot.png}
  \captionof{figure}{A helpful caption.}
\end{frame}
\end{verbatim}

Resize options: use \verb|width=|, \verb|height=|, \verb|scale=|, and \verb|keepaspectratio|.
\begin{verbatim}
\includegraphics[width=\linewidth]{img.png}
\includegraphics[height=0.4\textheight,keepaspectratio]{img.png}
\includegraphics[scale=0.8]{img.png}
\end{verbatim}

Two images side-by-side (columns):
\begin{verbatim}
\begin{frame}{Side-by-Side}
\begin{columns}[T]
  \begin{column}{0.48\textwidth}
    \centering
    \includegraphics[width=\linewidth]{img1.png}
    \captionof{figure}{Left}
  \end{column}
  \begin{column}{0.48\textwidth}
    \centering
    \includegraphics[width=\linewidth]{img2.png}
    \captionof{figure}{Right}
  \end{column}
\end{columns}
\end{frame}
\end{verbatim}

\section*{Tables}
\begin{verbatim}
\begin{frame}{Table}
\centering
\begin{tabular}{lcr}
\toprule
Left & Center & Right \\
\midrule
A    & B      & C \\
\bottomrule
\end{tabular}
\end{frame}
\end{verbatim}

\section*{Code Listings (listings)}
Global style is already set in \texttt{main.tex}. Per-snippet options:
\begin{verbatim}
\begin{lstlisting}[language=Python, caption={Quick demo}]
def f(x):
    return x**2
\end{lstlisting}
\end{verbatim}

\section*{Citations and References (biblatex)}
Add entries to \texttt{refs.bib}, then cite and print:
\begin{verbatim}
As shown in \cite{key2025}, ...
\begin{frame}[allowframebreaks]{References}
  \printbibliography
\end{frame}
\end{verbatim}

\section*{Including External Content}
Keep slides modular and include:
\begin{verbatim}
% In main.tex
\input{content/01_intro}
\input{content/02_theory}
\end{verbatim}


\end{document}
